\begin{table}[htbp]
\centering
\small
\caption{Network distance to the nearest facility, km (median [interquartile range]).}
\label{tab:t5_distances}
\begin{tabular}{lllll}
\toprule
City & Major road & Rail goods shed & IWT terminal & ICD \\
\midrule
Patna & 0.27 [0.06–0.65] & 13.14 [8.51–18.51] & 6.55 [4.18–11.50] & 32.43 [26.90–37.47] \\
Indore & 0.24 [0.07–0.56] & 4.61 [2.71–6.98] & 952$^{a}$ & 21.52 [18.91–25.00] \\
Varanasi & 0.49 [0.17–0.89] & 4.99 [3.27–7.65] & 5.32 [2.55–9.02] & 11.42 [8.77–16.08] \\
Kochi & 0.16 [0.02–0.47] & 6.11 [3.41–15.69] & 7.36 [4.51–16.17] & 11.09 [6.50–21.97] \\
Kolkata & 0.16 [0.04–0.40] & 5.69 [4.42–6.95] & 5.10 [2.41–7.45] & 5.99 [4.59–7.70] \\
\bottomrule
\end{tabular}
\begin{tablenotes}\footnotesize
\item $^{a}$ Out-of-district terminal, constant for all enterprises. The median enterprise in every city lies within 500 m of a major road; ICD distances vary more than fivefold across cities.
\end{tablenotes}
\end{table}
